\documentclass{article}
\usepackage{graphicx} % Required for inserting images
\usepackage{amsmath}
\usepackage{hyperref}
\usepackage{booktabs}

\title{Shopenstein: An AI-Vision-Based Shopping Navigator \\[0.5em] \large Embedded System Design}

\author{Group 4}
\date{October 2026}

\begin{document}

\maketitle

\section{Project name}

Shopenstein: An AI-Vision-Based Shopping Navigator.

\section{Project description}

This project implements an autonomous robotic navigator designed to traverse indoor shopping environments without relying on fixed floor line-tracking or physical tracks. The system integrates the Seeed Wio Terminal (ATSAMD51) as the master controller, the Grove Vision AI Module V2 (Himax HX6538) for onboard computer vision, and a Seeed Studio XIAO ESP32-C3 coprocessor for wireless communication. The vehicle detects predefined visual landmarks placed along shelves or walls, calculates relative spatial orientation and distance to them, and makes autonomous navigation and stopping decisions while enabling wireless teleoperation through a mobile web dashboard.

\section{Function description}

The AI-Vision-Based Shopping Navigator is a cyber-physical embedded system designed to automate and assist in indoor traversal. Using embedded computer vision at the edge, the vehicle detects visual landmark markers to localize itself, follow target waypoints, and provide visual diagnostics both on its onboard display and across a real-time wireless telemetry interface.

\subsection{Input}

The system processes three primary categories of inputs:

\begin{enumerate}
    \item \textbf{Visual Optical Input:} Real-time video captured by the onboard OV5647 CSI camera module connected directly to the Grove Vision AI Module V2.
    \item \textbf{Landmark Marker Signatures:} Trained visual classes representing specific landmark tags or objects, identified through edge neural network inference running directly on the vision coprocessor.
    \item \textbf{User Wireless Teleoperation Inputs:} Directional and proportional steering commands ($x, y$ coordinates) generated by a touch-based virtual joystick on a smartphone web browser, alongside mission-control commands such as \texttt{START AUTONOMOUS} and \texttt{STOP}.
\end{enumerate}

\subsection{Output}

The system generates the following outputs:

\begin{enumerate}
    \item \textbf{Differential Motor Actuation:} Two-channel Pulse Width Modulation (PWM) signals and H-bridge direction logic delivered to the TB6612FNG motor driver to control the two DC gear motors, incorporating software-based motor trimming to guarantee straight-line traversal.
    \item \textbf{Onboard Graphical Diagnostics:} A real-time diagnostic Heads-Up Display (HUD) on the Wio Terminal's built-in 320x240 TFT LCD showing vehicle state, detected marker ID, confidence percentage, bounding box coordinates, estimated distance, and neural network inference latency.
    \item \textbf{Wireless Web Dashboard and Radar Canvas:} Telemetry packets transmitted over serial to the ESP32-C3, which serves an interactive HTML5/Canvas radar and status monitor to the user's mobile device via Wi-Fi.
    \item \textbf{Autonomous Target Braking:} Automated closed-loop stopping action triggered when the target landmark reaches a calibrated proximity threshold.
\end{enumerate}

\subsection{How the system works}

The system operates as a distributed master-coprocessor embedded architecture composed of six interconnected functional modules:

\begin{enumerate}
    \item \textbf{User Interface (Smartphone Web Dashboard):} The user accesses a local web application served by the robot. The interface features a touch-sensitive virtual joystick that computes normalized horizontal ($x$) and vertical ($y$) deflections, state mode toggles, and an HTML5 Canvas rendering a real-time 2D vision radar.
    \item \textbf{Communications Coprocessor (Seeed XIAO ESP32-C3):} Configured as a standalone Wi-Fi Access Point (SoftAP) running an HTTP web server. It translates incoming HTTP touch events into serial commands (\texttt{JOY:x,y}, \texttt{CMD:AUTO}, \texttt{CMD:STOP}) and streams outgoing telemetry from the master MCU back to the client at 300~ms intervals.
    \item \textbf{Master Controller (Seeed Wio Terminal):} The core computational unit executing the primary finite-state machine (FSM). It communicates with the ESP32-C3 via UART (\texttt{Serial1} at 115200~bps) and with the vision sensor via Grove I2C. The Wio Terminal computes differential drive PWM values, manages the local LCD display, and executes autonomous navigation logic.
    \item \textbf{Edge Vision Sensor (Grove Vision AI Module V2):} Powered by the Himax HX6538 processor, this module captures frames from the OV5647 camera and executes an embedded Swift-YOLO object detection model. Upon receiving an \texttt{AT+INVOKE} command over I2C, it outputs bounding box coordinates ($x, y, w, h$), class target ID, and detection confidence.
    \item \textbf{Power and Driver Stage (Toshiba TB6612FNG):} A dual H-bridge motor driver that amplifies logic-level PWM signals from the Wio Terminal to deliver high-current drive to the motors.
    \item \textbf{Mechanical Actuation (2WD TT Gear Motors):} Two DC gear motors paired with rubber wheels and an omnidirectional front caster wheel provide differential steering and physical mobility.
\end{enumerate}

\begin{figure}[htbp]
    \centering
    % \includegraphics[width=0.9\textwidth]{figures/block_diagram.png}
    \fbox{\parbox[c][5.5cm][c]{0.88\textwidth}{\centering \textbf{[Figure Placeholder: System Functional Block Diagram (Inputs \& Outputs)]}}}
    \caption{Functional Block Diagram illustrating inputs, outputs, and communication channels between all system modules.}
    \label{fig:block_diagram}
\end{figure}

The system executes two complementary operational control loops:

\begin{itemize}
    \item \textbf{Manual Teleoperation Loop:} When the user deflects the virtual joystick, touch coordinates are pushed over Wi-Fi to the ESP32-C3 and forwarded via UART to the Wio Terminal. The Wio Terminal mixes throttle and steering into differential motor speeds, applies individual motor calibration factors (compensating for physical motor friction imbalances), and drives the TB6612FNG. Releasing the joystick automatically centers the controls and brings the robot to an immediate halt.
    \item \textbf{Autonomous Vision-Guided Loop:} Upon triggering the autonomous mode, the Wio Terminal commands the robot to drive forward at a stable base velocity while continuously polling the Grove Vision AI module over I2C. When the trained landmark is detected with a confidence score above the threshold ($Score \ge 50\%$), the master computes the estimated distance based on the marker's bounding box width:
    \begin{equation}
        d \approx \frac{f \cdot W}{w}
    \end{equation}
    where $W$ is the known marker width, $w$ is the detected width in pixels, and $f$ is the calibrated focal constant. When the marker size exceeds the arrival threshold ($w \ge 75\text{ px}$ or $d \le 25\text{ cm}$), the Wio Terminal cuts PWM drive to 0, transitions to the arrival state, displays a mission completion alert on its LCD, and updates the mobile interface.
\end{itemize}

\subsection{Circuit diagram}

To guarantee hardware safety and eliminate electromagnetic noise, the electrical system implements a dual-power architecture with a shared ground reference:
\begin{itemize}
    \item \textbf{Logic and Processing Power:} A 5V USB-C Power Bank powers the Wio Terminal. The Wio Terminal distributes 3.3V logic power to its onboard components and the TB6612FNG logic pins (\texttt{VCC} and \texttt{STBY}), while providing 5V (via pin 2 of its 40-pin header) to power the XIAO ESP32-C3.
    \item \textbf{Motor Power:} A separate 4x AA battery pack ($\sim$6V DC) connects exclusively to the \texttt{VM} power input of the TB6612FNG driver, ensuring that motor current surges do not induce brownouts or resets on the microcontrollers.
    \item \textbf{Common Ground:} The ground lines of the battery pack, motor driver, Wio Terminal, and ESP32-C3 are tied together to establish a common electrical reference.
\end{itemize}

\begin{figure}[htbp]
    \centering
    % \includegraphics[width=0.9\textwidth]{figures/circuit_diagram.png}
    \fbox{\parbox[c][5.5cm][c]{0.88\textwidth}{\centering \textbf{[Figure Placeholder: System Hardware Interconnection and Wiring Schematic]}}}
    \caption{Hardware wiring diagram showing power isolation, I2C bus, UART links, and motor driver pinout.}
    \label{fig:circuit_diagram}
\end{figure}

\section{Hardware used}

\begin{itemize}
    \item \textbf{Microcontroller (Master):} Seeed Wio Terminal (Microchip ATSAMD51P19 ARM Cortex-M4F core @ 120~MHz, built-in 320x240 SPI TFT LCD, 40-pin Raspberry Pi-compatible GPIO header).
    \item \textbf{Wireless Coprocessor:} Seeed Studio XIAO ESP32-C3 (RISC-V 32-bit single-core @ 160~MHz, 2.4~GHz Wi-Fi 802.11b/g/n).
    \item \textbf{AI Vision Module:} Grove Vision AI Module V2 (Himax WiseEye2 HX6538 dual-core Arm Cortex-M55 + Ethos-U55 NPU).
    \item \textbf{Camera Sensor:} Raspberry Pi Camera V1.3 compatible OV5647 5MP sensor connected via CSI ribbon cable.
    \item \textbf{Motor Driver:} Toshiba TB6612FNG Dual H-Bridge Driver module.
    \item \textbf{Chassis and Drive Train:} 2WD Smart Robot Car Chassis with two TT DC gear motors (1:48 reduction ratio) and one omnidirectional front caster wheel.
    \item \textbf{Power Sources:}
    \begin{itemize}
        \item 5V Portable Power Bank connected via USB-C (logic, displays, and sensors).
        \item 4x AA Battery Pack ($\sim$6V DC) dedicated to motor supply (\texttt{VM}).
    \end{itemize}
    \item \textbf{Interconnects and Prototyping:}
    \begin{itemize}
        \item 4-pin Grove cable (I2C bus connecting Wio Terminal and Vision AI Module).
        \item 400-point solderless breadboard.
        \item Assorted Male-to-Male and Male-to-Female Dupont jumper wires.
    \end{itemize}
    \item \textbf{Visual Targets:} High-contrast geometric landmark markers printed on matte paper.
\end{itemize}

\section{Software used}

\begin{itemize}
    \item \textbf{Firmware Development Toolchain:} Arduino IDE 2.x and \texttt{arduino-cli} on Linux.
    \item \textbf{Board Support Packages:}
    \begin{itemize}
        \item \texttt{Seeeduino:samd} (version 1.8.6) for the Wio Terminal.
        \item \texttt{esp32:esp32} (version 3.3.12) for the Seeed XIAO ESP32-C3.
    \end{itemize}
    \item \textbf{Embedded Libraries:}
    \begin{itemize}
        \item \texttt{Seeed\_Arduino\_SSCMA} (version 1.0.3) for packet-based I2C communication with the vision sensor.
        \item \texttt{TFT\_eSPI} / \texttt{Seeed\_Arduino\_LCD} for graphic rendering on the integrated LCD.
        \item \texttt{Wire} and \texttt{HardwareSerial} for I2C and UART communications.
        \item \texttt{WiFi} and \texttt{WebServer} for ESP32-C3 wireless Access Point and HTTP services.
    \end{itemize}
    \item \textbf{Edge Machine Learning Platform:} Seeed Studio SenseCraft AI for dataset management, model training, and deployment of Swift-YOLO to the Himax HX6538 NPU.
    \item \textbf{Web Dashboard Stack:} Handcrafted HTML5, CSS3, and JavaScript featuring vanilla touch event listeners for the virtual joystick and an HTML5 Canvas API for real-time 2D radar rendering.
\end{itemize}

\section{Reference}

% Placeholder for project references, datasheets, and academic literature
\begin{itemize}
    \item \textit{[Placeholder: Add project references, datasheets, textbooks, and related publications here]}
\end{itemize}

\end{document}
